% !TeX root= ../DoABook.tex

\chapter*{پیشگفتار}

\textit{تخمین جهت ورود سیگنال}\LTRfootnote{Direction-of-Arrival} که اغلب به‌اختصار \lr{DOA} نیز گفته می‌شود (یا جهت‌یابی)، به‌طور اساسی به برآورد جهت ورود سیگنال‌ها می‌پردازد؛ سیگنال‌هایی که به‌صورت امواج الکترومغناطیسی یا آکوستیکی به یک حسگر یا آرایه آنتن برخورد می‌کنند. نیاز به تخمین جهت ورود سیگنال از کاربردهایی مانند جست‌وجو و نجات، اجرای قانون، سونار، لرزه‌شناسی و مکان‌یابی تماس‌های اضطراری $911$ در شبکه‌های بی‌سیم ناشی می‌شود.

در حوزه پردازش سیگنال آرایه‌ای، نظریه‌ها و روش‌های متعددی برای تخمین جهت ورود سیگنال توسعه یافته‌اند و ادبیات گسترده‌ای نیز پیرامون آن وجود دارد. با وجود این، در جریان پژوهش‌ها و پیاده‌سازی‌های عملی دریافتیم که منابع موجود بسیار پراکنده‌اند؛ به‌طوری‌که برای مبتدیان یا دانشجویانی که می‌خواهند در مدت‌زمانی کوتاه وارد این حوزه شوند، دسترسی به یک منبع یکپارچه دشوار است. به بیان دیگر، کتاب‌های مروری اندکی وجود دارند که اصول و تکنیک‌های پایه تخمین جهت ورود سیگنال را به‌صورت منظم و زیر یک سقف ارائه کنند.

این کتاب با هدف رفع این نیاز تألیف شده است و مروری یکپارچه بر الگوریتم‌های پایه تخمین جهت ورود سیگنال، تحلیل کارایی آن‌ها و مقایسه میان روش‌ها ارائه می‌دهد. به‌طور ویژه، توصیف‌های نظام‌مند، تحلیل کارایی و مقایسه انواع الگوریتم‌ها ارائه شده و در پایان تمرکز خاصی بر خانواده روش \lr{ESPRIT} قرار گرفته است؛ روشی که نام آن از \textit{تخمین پارامترهای سیگنال از طریق تکنیک‌های ناوردایی دورانی}\LTRfootnote{Estimation of Signal Parameters via Rotational Invariance Techniques} گرفته شده است.

این کتاب برای مبتدیانی نظیر دانشجویان تحصیلات تکمیلی، مهندسان یا نهادهای قانون‌گذار که باید در مدت کوتاهی با مبانی تخمین جهت ورود سیگنال آشنا شوند، مناسب است. همچنین برای متخصصانی که قصد مرور دوباره مبانی این حوزه را دارند سودمند خواهد بود. امید است که این کتاب اطلاعات نظری کافی درباره مبانی تخمین جهت ورود سیگنال در اختیار خواننده قرار دهد تا بتواند مفاهیم اولیه را درک کرده و در صورت تمایل به موضوعات پیشرفته‌تر وارد شود.






